\documentclass[aspectratio=169,10pt]{beamer}

% ---------- Packages ----------
\usepackage[utf8]{inputenc}
\usepackage[T1]{fontenc}
\usepackage{amsmath,amssymb,amsthm}
\usepackage{booktabs}
\usepackage{xcolor}

% ---------- Theme ----------
\usetheme{Madrid}
\usecolortheme{seahorse}
\setbeamertemplate{navigation symbols}{}
\setbeamertemplate{footline}[frame number]
\definecolor{ans}{RGB}{0,70,140}

% ---------- Helper macro: \qa{label}{question}{worked answer} ----------
\newcommand{\qa}[3]{%
  \textbf{Q#1.} #2\par\smallskip
  {\color{ans}\textbf{Solution.} #3}\par\vspace{6pt}}

% ---------- Title info ----------
\title{Mathematics of Computing}
\subtitle{Module 1 Assignment: Logic, Proofs and Induction}
\author{Aishwarya J S}
\institute{M.Tech CSE (Artificial Intelligence and Data Science)}
\date{\today}

\begin{document}

\begin{frame}
  \titlepage
\end{frame}

\begin{frame}{Outline}
\small
\begin{columns}[t]
\begin{column}{0.48\textwidth}
\begin{enumerate}
  \item Proposition
  \item Predicate
  \item Quantifiers
  \item Truth Table
  \item First Order Logic
  \item Satisfiability
  \item Patterns of Proof
\end{enumerate}
\end{column}
\begin{column}{0.48\textwidth}
\begin{enumerate}\setcounter{enumi}{7}
  \item Proving If and Only If
  \item Proving Sets and Set Equations
  \item Russell's Paradox
  \item Well Ordering Principle
  \item Induction and Invariants
  \item Strong Induction
  \item Structural Induction
\end{enumerate}
\end{column}
\end{columns}
\medskip
\centering Each topic: 5 constructive questions with worked solutions.
\end{frame}

% ================= 1. Proposition =================
\begin{frame}{1. Proposition (Q1--Q2)}
\footnotesize
\qa{1}{Which are propositions? (a) $x+2=5$ (b) Is 7 prime? (c) ``This statement is false.'' (d) Every even integer $>2$ is a sum of two primes. (e) $2+3=5$.}
{A proposition is a declarative sentence that is exactly one of true or false.\\
(a) Its truth depends on the value of $x$, so it is a predicate, not a proposition.\\
(b) A question asserts nothing, so it has no truth value.\\
(c) If it were true it would be false, and if false it would be true, so no consistent value exists.\\
(d) This is Goldbach's conjecture. Nobody knows its value, but it is definitely true or false.\\
(e) True.\\
\emph{Hence only (d) and (e) are propositions.}}
\qa{2}{Negate ``If it rains, the match is cancelled.''}
{Let $p$ = ``it rains'' and $q$ = ``the match is cancelled''. The sentence is $p\to q\equiv\neg p\lor q$.\\
Negating: $\neg(\neg p\lor q)\equiv p\land\neg q$ (De Morgan).\\
\emph{Hence the negation is: ``It rains and the match is not cancelled.''} It is the only situation that breaks the promise.}
\end{frame}

\begin{frame}{1. Proposition (Q3--Q4)}
\footnotesize
\qa{3}{Translate ``You can access the site only if you are registered and verified.''}
{Let $a$ = access, $r$ = registered, $v$ = verified.\\
``$A$ only if $B$'' says $B$ is necessary for $A$: whenever $A$ holds, $B$ holds. So it is $A\to B$, not $B\to A$.\\
Here $B$ is ``$r$ and $v$'', so the statement is $a\to(r\land v)$.\\
\emph{Check:} a registered but unverified user cannot access, since $a$ would make $r\land v$ false.}
\qa{4}{Build a compound proposition in $p,q,r$ that is true exactly when exactly one of them is true.}
{List the rows where exactly one variable is T: $(1,0,0)$, $(0,1,0)$, $(0,0,1)$.\\
Write one conjunction per row, then join with $\lor$:\\
\centerline{$(p\land\neg q\land\neg r)\lor(\neg p\land q\land\neg r)\lor(\neg p\land\neg q\land r)$}
Each disjunct is true on exactly one row, so the whole formula is true on exactly those three rows.}
\end{frame}

\begin{frame}{1. Proposition (Q5)}
\footnotesize
\qa{5}{For ``If $n$ is divisible by 6, then $n$ is divisible by 3'' write the converse, inverse and contrapositive. Which are equivalent to the original?}
{Let $p$ = ``$6\mid n$'' and $q$ = ``$3\mid n$''. The original $p\to q$ is true, since $n=6k=3(2k)$.
\begin{itemize}\footnotesize
 \item Converse $q\to p$: if $3\mid n$ then $6\mid n$. False: $n=3$.
 \item Inverse $\neg p\to\neg q$: if $6\nmid n$ then $3\nmid n$. False: $n=3$.
 \item Contrapositive $\neg q\to\neg p$: if $3\nmid n$ then $6\nmid n$. True: if $6\mid n$ then $3\mid n$.
\end{itemize}
\emph{Hence only the contrapositive is equivalent.} The converse and inverse are equivalent to each other, but not to the original.}
\end{frame}

% ================= 2. Predicate =================
\begin{frame}{2. Predicate (Q1--Q2)}
\footnotesize
\qa{1}{Let $P(x)$: ``$x$ is prime'' and $E(x)$: ``$x$ is even'' over the integers. Express ``2 is the only even prime.''}
{The claim has two parts.\\
\emph{Existence:} 2 is an even prime, i.e.\ $P(2)\land E(2)$.\\
\emph{Uniqueness:} any integer that is both prime and even must be 2, i.e.\ $\forall x\,((P(x)\land E(x))\to x=2)$.\\
\emph{Hence:} $P(2)\land E(2)\land\forall x\,((P(x)\land E(x))\to x=2)$.\\
Without the second part the formula would not rule out other even primes.}
\qa{2}{Over $\mathbb{R}$ let $P(x,y)$: $x+y=0$. Decide the truth of $\forall x\exists y\,P(x,y)$ and $\exists y\forall x\,P(x,y)$.}
{$\forall x\exists y$: $y$ may depend on $x$. Given any $x$, choose $y=-x$; then $x+y=0$. \textbf{True.}\\
$\exists y\forall x$: one $y$ must work for every $x$. Suppose $y_0$ works. Then $x=1-y_0$ gives $x+y_0=1\neq0$. \textbf{False.}\\
\emph{Hence swapping quantifiers changes meaning.}}
\end{frame}

\begin{frame}{2. Predicate (Q3--Q4)}
\footnotesize
\qa{3}{Find the truth set of $P(x)$: $x^2>x$ over $\mathbb{N}$ and over $\mathbb{R}$.}
{Rewrite as $x^2-x>0$, i.e.\ $x(x-1)>0$. A product is positive when both factors have the same sign.\\
Both positive: $x>0$ and $x>1$, so $x>1$. Both negative: $x<0$ and $x<1$, so $x<0$.\\
Over $\mathbb{R}$ the truth set is $(-\infty,0)\cup(1,\infty)$.\\
Over $\mathbb{N}=\{0,1,2,\dots\}$ no negatives exist, and $x=0,1$ give equality, so the set is $\{2,3,4,\dots\}$.}
\qa{4}{Express ``No student failed every exam'' and push the negation inward.}
{With $F(s,e)$: ``student $s$ failed exam $e$'', the sentence is $\neg\exists s\,\forall e\,F(s,e)$.\\
Step 1: $\neg\exists s\,(\cdot)\equiv\forall s\,\neg(\cdot)$, giving $\forall s\,\neg\forall e\,F(s,e)$.\\
Step 2: $\neg\forall e\,(\cdot)\equiv\exists e\,\neg(\cdot)$, giving $\forall s\,\exists e\,\neg F(s,e)$.\\
\emph{Reading back:} every student passed at least one exam.}
\end{frame}

\begin{frame}{2. Predicate (Q5)}
\footnotesize
\qa{5}{Over positive integers let $D(x,y)$: ``$x$ divides $y$''. Find all $x$ with $D(x,12)\land D(x,18)$.}
{Divisors of 12: $\{1,2,3,4,6,12\}$. Divisors of 18: $\{1,2,3,6,9,18\}$.\\
The truth set of $D(x,12)\land D(x,18)$ is the intersection: $\{1,2,3,6\}$.\\
\emph{Cross-check:} $\gcd(12,18)=6$, and the common divisors of two numbers are exactly the divisors of their gcd, namely $1,2,3,6$.}
\end{frame}

% ================= 3. Quantifiers =================
\begin{frame}{3. Quantifiers (Q1--Q2)}
\footnotesize
\qa{1}{Negate the $\varepsilon$--$\delta$ statement of continuity at $a$: $\forall\varepsilon>0\,\exists\delta>0\,\forall x\,(|x-a|<\delta\to|f(x)-f(a)|<\varepsilon)$.}
{Move the negation across one symbol at a time.\\
$\forall\varepsilon\to\exists\varepsilon$, then $\exists\delta\to\forall\delta$, then $\forall x\to\exists x$.\\
Inside, $\neg(p\to q)\equiv p\land\neg q$.\\
\emph{Result:} $\exists\varepsilon>0\,\forall\delta>0\,\exists x\,(|x-a|<\delta\land|f(x)-f(a)|\ge\varepsilon)$.\\
In words: some tolerance $\varepsilon$ exists such that, however small $\delta$ is, an $x$ within $\delta$ of $a$ has $f(x)$ at least $\varepsilon$ away from $f(a)$.}
\qa{2}{Is $\forall x\exists y\,(y^2=x)$ true over $\mathbb{R}$? Over $\mathbb{R}_{\ge0}$?}
{Over $\mathbb{R}$: take $x=-1$. No real $y$ has $y^2=-1$, so the statement is \textbf{false}.\\
Over $\mathbb{R}_{\ge0}$: for any $x\ge0$ choose $y=\sqrt x$, which is real. \textbf{True.}\\
\emph{The domain decides the truth value.}}
\end{frame}

\begin{frame}{3. Quantifiers (Q3--Q4)}
\footnotesize
\qa{3}{Show that $\forall x(P(x)\lor Q(x))$ is not equivalent to $\forall x\,P(x)\lor\forall x\,Q(x)$.}
{It is enough to give one structure where they differ. Take domain $\mathbb{Z}$, $P(x)$ = ``$x$ is even'', $Q(x)$ = ``$x$ is odd''.\\
Left side: every integer is even or odd, so it is \textbf{true}.\\
Right side: ``all integers are even'' is false ($1$), and ``all are odd'' is false ($2$), so it is \textbf{false}.\\
\emph{Hence $\forall$ does not distribute over $\lor$.} (By contrast $\forall x(P\land Q)\equiv\forall xP\land\forall xQ$.)}
\qa{4}{Express ``exactly two elements satisfy $P$.''}
{We need ``at least two'' and ``no third''.\\
At least two: $\exists x\exists y\,(x\ne y\land P(x)\land P(y))$.\\
No others: $\forall z\,(P(z)\to(z=x\lor z=y))$.\\
\emph{Combined:} $\exists x\exists y\,[x\ne y\land P(x)\land P(y)\land\forall z\,(P(z)\to(z=x\lor z=y))]$.}
\end{frame}

\begin{frame}{3. Quantifiers (Q5)}
\footnotesize
\qa{5}{Define $\exists!x\,P(x)$ using only $\exists$, $\forall$, $=$. Is $\exists!x\,(x^2=4)$ true over $\mathbb{R}$? Over $\mathbb{R}^{+}$?}
{``Exactly one'' means a witness exists and any other witness equals it:\\
\centerline{$\exists!x\,P(x)\equiv\exists x\,[P(x)\land\forall y\,(P(y)\to y=x)]$.}
Over $\mathbb{R}$ the solutions of $x^2=4$ are $2$ and $-2$, so uniqueness fails: \textbf{false}.\\
Over $\mathbb{R}^{+}$ only $2$ qualifies, so \textbf{true}.}
\end{frame}

% ================= 4. Truth Table =================
\begin{frame}{4. Truth Table (Q1)}
\footnotesize
\textbf{Q1.} Show $((p\to q)\land(q\to r))\to(p\to r)$ is a tautology.\par\smallskip
\centering
\begin{tabular}{ccc|c|c|c|c|c}
\toprule
$p$&$q$&$r$&$p\to q$&$q\to r$&$(p\to q)\land(q\to r)$&$p\to r$&Whole\\
\midrule
T&T&T&T&T&T&T&T\\
T&T&F&T&F&F&F&T\\
T&F&T&F&T&F&T&T\\
T&F&F&F&T&F&F&T\\
F&T&T&T&T&T&T&T\\
F&T&F&T&F&F&T&T\\
F&F&T&T&T&T&T&T\\
F&F&F&T&T&T&T&T\\
\bottomrule
\end{tabular}\par\smallskip
{\color{ans}\raggedright\textbf{Solution.} With 3 variables there are $2^3=8$ rows. The last column is T in every row, so the formula is a tautology (hypothetical syllogism). The only rows where the antecedent is T are rows 1, 5, 7, 8, and there $p\to r$ is T too.\par}
\end{frame}

\begin{frame}{4. Truth Table (Q2--Q3)}
\footnotesize
\qa{2}{Verify $p\oplus q\equiv(p\lor q)\land\neg(p\land q)$.}
{\centering
\begin{tabular}{cc|c|c|c|c}
\toprule
$p$&$q$&$p\oplus q$&$p\lor q$&$\neg(p\land q)$&RHS\\
\midrule
T&T&F&T&F&F\\
T&F&T&T&T&T\\
F&T&T&T&T&T\\
F&F&F&F&T&F\\
\bottomrule
\end{tabular}\par\smallskip
\raggedright The columns $p\oplus q$ and RHS agree in all four rows, so the two formulas are equivalent.}
\qa{3}{For a formula with $n$ variables, how many rows does its truth table have, and how many distinct formulas (up to equivalence) exist? Evaluate for $n=3$.}
{Each variable is T or F independently, so there are $2^n$ rows. A formula is determined by its output column; each of the $2^n$ rows can output T or F independently, so there are $2^{2^n}$ possible columns.\\
For $n=3$: $8$ rows and $2^8=256$ distinct functions.}
\end{frame}

\begin{frame}{4. Truth Table (Q4--Q5)}
\footnotesize
\qa{4}{Decide whether $(p\lor q)\to r\equiv(p\to r)\land(q\to r)$.}
{\emph{Algebraic route.} $(p\lor q)\to r\equiv\neg(p\lor q)\lor r\equiv(\neg p\land\neg q)\lor r$. By distributivity this is $(\neg p\lor r)\land(\neg q\lor r)\equiv(p\to r)\land(q\to r)$.\\
\emph{Truth-table check.} The left side is false only when $r=\mathrm{F}$ and ($p$ or $q$) is T. The right side is false only when $r=\mathrm{F}$ and ($p=$T or $q=$T). The falsifying rows coincide.\\
\emph{Hence they are equivalent.}}
\qa{5}{Construct formulas for: (a) true only on rows TT and FF of $(p,q)$; (b) the 3-variable majority function.}
{(a) Write one conjunction for each true row and join them: $(p\land q)\lor(\neg p\land\neg q)$, which is $p\leftrightarrow q$.\\
(b) Majority is true when at least two of $p,q,r$ are T. Each pair gives a disjunct: $(p\land q)\lor(p\land r)\lor(q\land r)$. This covers TTT, TTF, TFT, FTT and no other row.}
\end{frame}

% ================= 5. First Order Logic =================
\begin{frame}{5. First Order Logic (Q1--Q2)}
\footnotesize
\qa{1}{Using $\mathrm{Student}(x)$, $\mathrm{Course}(y)$, $\mathrm{Takes}(x,y)$, write ``Some student takes every course.''}
{We need one student $x$ such that for all $y$, if $y$ is a course then $x$ takes $y$.\\
\centerline{$\exists x\,[\mathrm{Student}(x)\land\forall y\,(\mathrm{Course}(y)\to\mathrm{Takes}(x,y))]$}
The $\exists$ is paired with $\land$ and the $\forall$ with $\to$. Swapping them would change the meaning.}
\qa{2}{Give a structure where $\forall x\exists y\,R(x,y)$ is true but $\exists y\forall x\,R(x,y)$ is false.}
{Take domain $\mathbb{N}$ and $R(x,y)$ interpreted as $x<y$.\\
$\forall x\exists y$: for any $x$, choose $y=x+1$. True.\\
$\exists y\forall x$: a single $y$ would have to exceed every $x$, including $x=y$, and $y<y$ is false. False.}
\end{frame}

\begin{frame}{5. First Order Logic (Q3--Q4)}
\footnotesize
\qa{3}{Express ``$f$ is injective'' and ``$f$ is surjective''. Test $f(n)=2n$ on $\mathbb{Z}$.}
{Injective: $\forall x\forall y\,(f(x)=f(y)\to x=y)$. Surjective: $\forall y\exists x\,f(x)=y$.\\
Injective test: $2x=2y$ gives $x=y$, so yes.\\
Surjective test: for $y=1$ we would need $2x=1$, impossible in $\mathbb{Z}$, so no.\\
\emph{Hence $f$ is injective but not surjective.}}
\qa{4}{Show $\forall x\,P(x)\to\exists x\,P(x)$ is valid in first order logic.}
{Validity means true in every structure. By the usual convention a structure's domain is nonempty. Suppose $\forall x\,P(x)$ holds in some structure, and pick any element $a$ of its domain. Then $P(a)$ holds, so $\exists x\,P(x)$ holds. If the antecedent fails, the implication is true anyway.\\
\emph{Hence it is valid.} (With empty domains allowed it would fail.)}
\end{frame}

\begin{frame}{5. First Order Logic (Q5)}
\footnotesize
\qa{5}{Show $\{\forall x(P(x)\to Q(x)),\ \exists x\,P(x)\}\models\exists x\,Q(x)$.}
{Let a structure make both premises true.\\
Step 1: by $\exists x\,P(x)$ there is an element $a$ with $P(a)$.\\
Step 2: instantiate $\forall x(P(x)\to Q(x))$ at $a$ to get $P(a)\to Q(a)$.\\
Step 3: modus ponens gives $Q(a)$.\\
Step 4: so some element satisfies $Q$, i.e.\ $\exists x\,Q(x)$.\\
\emph{Hence the conclusion holds in every structure satisfying the premises.}}
\end{frame}

% ================= 6. Satisfiability =================
\begin{frame}{6. Satisfiability (Q1--Q2)}
\footnotesize
\qa{1}{Is $(p\lor q)\land(\neg p\lor q)\land(p\lor\neg q)\land(\neg p\lor\neg q)$ satisfiable?}
{Each clause is false for exactly one assignment of $(p,q)$:\\
$p\lor q$ is false at FF; $\neg p\lor q$ at TF; $p\lor\neg q$ at FT; $\neg p\lor\neg q$ at TT.\\
The four clauses cover all four assignments, so every assignment falsifies one clause.\\
\emph{Hence the formula is unsatisfiable.}}
\qa{2}{Find a satisfying assignment for $(p\lor\neg q)\land(q\lor r)\land(\neg p\lor\neg r)$.}
{Try $p=\mathrm{T}$. The third clause forces $r=\mathrm{F}$, and then the second forces $q=\mathrm{T}$. Check: clause 1 is T$\lor$F $=$T, clause 2 is T$\lor$F $=$T, clause 3 is F$\lor$T $=$T.\\
\emph{Hence $p=$T, $q=$T, $r=$F satisfies it.}}
\end{frame}

\begin{frame}{6. Satisfiability (Q3--Q4)}
\footnotesize
\qa{3}{Use ``$A$ is valid iff $\neg A$ is unsatisfiable'' to show $((p\to q)\land p)\to q$ is valid.}
{Negate: $\neg(((p\to q)\land p)\to q)\equiv(p\to q)\land p\land\neg q$.\\
Suppose some assignment satisfies it. Then $p$ is T and $p\to q$ is T, so $q$ is T. But $\neg q$ is also T, a contradiction.\\
So the negation is unsatisfiable. \emph{Hence the original is valid (modus ponens).}}
\qa{4}{Show that $\forall x\exists y\,R(x,y)\land\forall x\,\neg R(x,x)\land(R\text{ transitive})$ is satisfiable but has no finite model.}
{\emph{Satisfiable:} $(\mathbb{N},<)$ works: every number has a larger one, $x<x$ never holds, and $<$ is transitive.\\
\emph{No finite model:} suppose one existed. Start at $a_0$ and pick $a_{i+1}$ with $R(a_i,a_{i+1})$. By transitivity $R(a_i,a_j)$ for all $i<j$. A finite domain forces a repeat $a_i=a_j$ with $i<j$, giving $R(a,a)$, contradicting irreflexivity. So every model is infinite.}
\end{frame}

\begin{frame}{6. Satisfiability (Q5)}
\footnotesize
\qa{5}{Encode ``2-colour the triangle $a,b,c$ with colours T/F'' as clauses. Is it satisfiable?}
{Let each variable be the colour of a vertex. An edge $(u,v)$ needs different colours: $(u\lor v)\land(\neg u\lor\neg v)$.\\
Edges: $(a\lor b)\land(\neg a\lor\neg b)$, $(b\lor c)\land(\neg b\lor\neg c)$, $(a\lor c)\land(\neg a\lor\neg c)$.\\
From the first two edges, $a\ne b$ and $b\ne c$, so with only two colours $a=c$. The third edge demands $a\ne c$, a contradiction.\\
\emph{Hence the formula is unsatisfiable: a triangle is not 2-colourable.}}
\end{frame}

% ================= 7. Patterns of Proof =================
\begin{frame}{7. Patterns of Proof (Q1--Q2)}
\footnotesize
\qa{1 (cases)}{Prove $n^2+n$ is even for every integer $n$.}
{Every integer is even or odd, so use two cases.\\
\emph{Case 1:} $n=2k$. Then $n^2+n=4k^2+2k=2(2k^2+k)$.\\
\emph{Case 2:} $n=2k+1$. Then $n^2+n=4k^2+4k+1+2k+1=2(2k^2+3k+1)$.\\
Both are even. \emph{Hence $n^2+n$ is even for all $n$.}}
\qa{2 (cases)}{Prove $3\mid n^3-n$ for every integer $n$.}
{Factor: $n^3-n=(n-1)\,n\,(n+1)$. Split by the remainder of $n$ modulo 3.\\
\emph{Case $n=3k$:} the factor $n$ is divisible by 3.\\
\emph{Case $n=3k+1$:} the factor $n-1=3k$ is divisible by 3.\\
\emph{Case $n=3k+2$:} the factor $n+1=3(k+1)$ is divisible by 3.\\
\emph{Hence in every case the product is divisible by 3.}}
\end{frame}

\begin{frame}{7. Patterns of Proof (Q3--Q4)}
\footnotesize
\qa{3 (direct proof of an implication)}{Prove: if $a\mid b$ and $b\mid c$ then $a\mid c$.}
{Assume $a\mid b$ and $b\mid c$. By definition there are integers $k,m$ with $b=ak$ and $c=bm$.\\
Substitute: $c=(ak)m=a(km)$, and $km$ is an integer.\\
\emph{Hence $a\mid c$.}}
\qa{4 (contrapositive)}{Prove: if $n^2$ is even then $n$ is even.}
{A direct proof is awkward, so prove the contrapositive: if $n$ is odd then $n^2$ is odd.\\
Let $n=2k+1$. Then $n^2=4k^2+4k+1=2(2k^2+2k)+1$, which is odd.\\
\emph{Since an implication equals its contrapositive, the original statement is proved.}}
\end{frame}

\begin{frame}{7. Patterns of Proof (Q5)}
\footnotesize
\qa{5 (contradiction)}{Prove that the sum of a rational $q$ and an irrational $x$ is irrational.}
{Assume the opposite: $q+x=r$ for some rational $r$.\\
Then $x=r-q$. The difference of two rationals is rational, so $x$ would be rational.\\
This contradicts the hypothesis that $x$ is irrational.\\
\emph{Hence $q+x$ is irrational.}}
\end{frame}

% ================= 8. Iff =================
\begin{frame}{8. Proving If and Only If (Q1--Q2)}
\footnotesize
\qa{1}{For integers $a,b$ prove: $a+b$ is even $\iff$ $a$ and $b$ have the same parity.}
{($\Leftarrow$) If both are even, $a=2k$, $b=2m$ and $a+b=2(k+m)$. If both are odd, $a=2k+1$, $b=2m+1$ and $a+b=2(k+m+1)$. Either way the sum is even.\\
($\Rightarrow$) Contrapositive: suppose the parities differ, say $a=2k$, $b=2m+1$. Then $a+b=2(k+m)+1$ is odd.\\
\emph{Both directions hold, so the equivalence is proved.}}
\qa{2}{For nonzero integers prove: ($a\mid b$ and $b\mid a$) $\iff$ $a=\pm b$.}
{($\Rightarrow$) Write $b=ak$ and $a=bm$. Then $a=akm$, and since $a\ne0$, $km=1$. Integers with product 1 are $\pm1$, so $k=\pm1$ and $b=\pm a$.\\
($\Leftarrow$) If $a=\pm b$ then $b=(\pm1)a$ and $a=(\pm1)b$, so each divides the other.}
\end{frame}

\begin{frame}{8. Proving If and Only If (Q3--Q4)}
\footnotesize
\qa{3}{For real $x$ prove $|x|<3\iff-3<x<3$.}
{\emph{Case $x\ge0$:} $|x|=x$, so $|x|<3\iff x<3$; and $x\ge0>-3$ already, so this is $-3<x<3$.\\
\emph{Case $x<0$:} $|x|=-x$, so $|x|<3\iff-x<3\iff x>-3$; and $x<0<3$ already, so this is $-3<x<3$.\\
Every step in both cases is reversible, so the equivalence holds.}
\qa{4}{For $a,b>0$ prove $a<b\iff a^2<b^2$.}
{($\Rightarrow$) $b^2-a^2=(b-a)(b+a)$. Here $b-a>0$ and $b+a>0$, so the product is positive.\\
($\Leftarrow$) Contrapositive: suppose $a\ge b$. Since both are positive, $a^2\ge ab\ge b^2$, so $a^2\ge b^2$, i.e.\ $a^2<b^2$ fails.}
\end{frame}

\begin{frame}{8. Proving If and Only If (Q5)}
\footnotesize
\qa{5}{Show $p\leftrightarrow q\equiv\neg p\leftrightarrow\neg q$.}
{$p\leftrightarrow q$ is true exactly when $p$ and $q$ have the same truth value.\\
Negating both variables turns T into F and F into T on both sides, so equal values stay equal and different values stay different.\\
\emph{Truth table:} on $(p,q)=$ TT, TF, FT, FF, the left side is T, F, F, T and the right side (with inputs FF, FT, TF, TT) is T, F, F, T. The columns match.}
\end{frame}

% ================= 9. Sets =================
\begin{frame}{9. Proving Sets and Set Equations (Q1--Q2)}
\footnotesize
\qa{1}{Prove $A\cap(B\cup C)=(A\cap B)\cup(A\cap C)$ by the element method.}
{Take an arbitrary $x$ and chain equivalences:\\
$x\in A\cap(B\cup C)\iff x\in A\land(x\in B\lor x\in C)$\\
$\iff(x\in A\land x\in B)\lor(x\in A\land x\in C)$ (distributivity of $\land$ over $\lor$)\\
$\iff x\in(A\cap B)\cup(A\cap C)$.\\
Both sets have exactly the same elements, so they are equal.}
\qa{2}{Prove $A\setminus(B\cap C)=(A\setminus B)\cup(A\setminus C)$ by set algebra.}
{Use $X\setminus Y=X\cap Y^c$.\\
$A\setminus(B\cap C)=A\cap(B\cap C)^c=A\cap(B^c\cup C^c)$ (De Morgan)\\
$=(A\cap B^c)\cup(A\cap C^c)$ (distributivity) $=(A\setminus B)\cup(A\setminus C)$.}
\end{frame}

\begin{frame}{9. Proving Sets and Set Equations (Q3--Q4)}
\footnotesize
\qa{3}{Prove $A\subseteq B\iff A\cup B=B$.}
{($\Rightarrow$) Assume $A\subseteq B$. Always $B\subseteq A\cup B$. Conversely, let $x\in A\cup B$: if $x\in A$ then $x\in B$ by assumption, otherwise $x\in B$ directly. So $A\cup B=B$.\\
($\Leftarrow$) Assume $A\cup B=B$. If $x\in A$ then $x\in A\cup B=B$. So $A\subseteq B$.}
\qa{4}{Prove or disprove $(A\cup B)\setminus C=A\cup(B\setminus C)$.}
{Try small sets: $A=\{1\}$, $B=\varnothing$, $C=\{1\}$.\\
Left: $(\{1\}\cup\varnothing)\setminus\{1\}=\varnothing$. Right: $\{1\}\cup(\varnothing\setminus\{1\})=\{1\}$.\\
\emph{They differ, so the identity is false.} The correct version is $(A\cup B)\setminus C=(A\setminus C)\cup(B\setminus C)$.}
\end{frame}

\begin{frame}{9. Proving Sets and Set Equations (Q5)}
\footnotesize
\qa{5}{Prove $\mathcal{P}(A)\cap\mathcal{P}(B)=\mathcal{P}(A\cap B)$. Does $\mathcal{P}(A)\cup\mathcal{P}(B)=\mathcal{P}(A\cup B)$?}
{\emph{Intersection:} $X\in\mathcal P(A)\cap\mathcal P(B)\iff X\subseteq A\land X\subseteq B\iff X\subseteq A\cap B\iff X\in\mathcal P(A\cap B)$.\\
\emph{Union:} the inclusion $\subseteq$ always holds, but equality fails. Take $A=\{1\}$, $B=\{2\}$. Then $\{1,2\}\in\mathcal P(A\cup B)$, but $\{1,2\}\not\subseteq A$ and $\{1,2\}\not\subseteq B$, so it is in neither $\mathcal P(A)$ nor $\mathcal P(B)$.}
\end{frame}

% ================= 10. Russell =================
\begin{frame}{10. Russell's Paradox (Q1--Q2)}
\footnotesize
\qa{1}{Derive Russell's paradox.}
{Naive comprehension allows $R=\{x\mid x\notin x\}$. Ask whether $R\in R$.\\
If $R\in R$, then $R$ must satisfy the defining condition, so $R\notin R$.\\
If $R\notin R$, then $R$ satisfies the condition, so $R\in R$.\\
Both cases give a contradiction. \emph{Hence unrestricted set formation is inconsistent.}}
\qa{2}{(Cantor) Show there is no surjection $f:A\to\mathcal P(A)$, using a Russell-style set.}
{Suppose $f$ is onto. Define $D=\{a\in A\mid a\notin f(a)\}\subseteq A$. Then $D\in\mathcal P(A)$, so $D=f(d)$ for some $d\in A$.\\
Is $d\in D$? $d\in D\iff d\notin f(d)=D$. So $d\in D\iff d\notin D$, a contradiction.\\
\emph{Hence no surjection exists, and $\mathcal P(A)$ is strictly larger than $A$.}}
\end{frame}

\begin{frame}{10. Russell's Paradox (Q3--Q4)}
\footnotesize
\qa{3}{How does the Separation axiom avoid the paradox?}
{Separation only allows $\{x\in A\mid\varphi(x)\}$ where $A$ is already a set.\\
Form $R_A=\{x\in A\mid x\notin x\}$. If $R_A\in A$, then $R_A\in R_A\iff R_A\notin R_A$, a contradiction. So $R_A\notin A$.\\
\emph{Hence the construction just shows that $R_A$ is a set outside $A$; no paradox arises.}}
\qa{4}{Prove there is no set of all sets.}
{Suppose $V$ is a set containing every set. By Separation, $R=\{x\in V\mid x\notin x\}$ is a set, so $R\in V$.\\
But the previous answer showed $R\notin V$. Contradiction.\\
\emph{Hence no universal set exists.}}
\end{frame}

\begin{frame}{10. Russell's Paradox (Q5)}
\footnotesize
\qa{5}{Show the Grelling--Nelson paradox has the same structure: call a word ``heterological'' if it does not describe itself.}
{``Short'' is a short word, so it describes itself (autological). ``Long'' is not long, so it is heterological.\\
Is ``heterological'' heterological?\\
If yes, it does not describe itself, so it is not heterological. If no, it describes itself, so it \emph{is} heterological.\\
This is Russell's argument with ``word describes itself'' in place of $x\in x$.}
\end{frame}

% ================= 11. WOP =================
\begin{frame}{11. Well Ordering Principle (Q1--Q2)}
\footnotesize
\qa{1}{State the principle. Does it hold for $\mathbb{Z}$, $\mathbb{Q}_{\ge0}$ and $[0,1]$?}
{\emph{WOP:} every nonempty subset of $\mathbb{N}$ has a least element.\\
$\mathbb{Z}$: the subset of negative integers has no least element.\\
$\mathbb{Q}_{\ge0}$: the positive rationals have no minimum, since $q/2$ is always smaller.\\
$[0,1]$: the subset $(0,1]$ has no minimum for the same reason.\\
\emph{Hence WOP is special to $\mathbb{N}$ (and sets order-isomorphic to it).}}
\qa{2}{Use WOP to prove every positive integer is a sum of distinct powers of 2.}
{Suppose not, and let $n$ be the least counterexample. $n\ne1=2^0$.\\
Let $2^k$ be the largest power of 2 with $2^k\le n$. If $n=2^k$ we are done, so $n>2^k$. Put $m=n-2^k$. Then $0<m<2^k$ (since $n<2^{k+1}$) and $m<n$.\\
By minimality of $n$, $m$ is a sum of distinct powers of 2, all at most $m<2^k$. Adding $2^k$ gives a representation of $n$ with distinct powers. Contradiction.}
\end{frame}

\begin{frame}{11. Well Ordering Principle (Q3--Q4)}
\footnotesize
\qa{3}{Use WOP to show every integer $n>1$ has a prime divisor.}
{Let $S$ be the set of divisors of $n$ that exceed 1. It is nonempty because $n\in S$. Let $d$ be its least element.\\
Suppose $d$ is not prime. Then $d=ab$ with $1<a<d$. Since $a\mid d$ and $d\mid n$, also $a\mid n$, so $a\in S$ and $a<d$, contradicting minimality.\\
\emph{Hence $d$ is a prime divisor of $n$.}}
\qa{4}{Use WOP to prove $1+2+\cdots+n=\frac{n(n+1)}{2}$ for all $n\ge1$.}
{Suppose the set of counterexamples is nonempty and let $m$ be its least element. $m\ne1$ because $1=\frac{1\cdot2}{2}$.\\
So $m-1\ge1$ is not a counterexample: $1+\cdots+(m-1)=\frac{(m-1)m}{2}$. Adding $m$: $\frac{m(m-1)}{2}+m=\frac{m(m+1)}{2}$.\\
So $m$ satisfies the formula, a contradiction.}
\end{frame}

\begin{frame}{11. Well Ordering Principle (Q5)}
\footnotesize
\qa{5}{Use WOP to prove the division algorithm: for $a\in\mathbb Z$ and $b>0$ there exist $q,r$ with $a=bq+r$ and $0\le r<b$.}
{Let $S=\{a-bk\mid k\in\mathbb Z,\ a-bk\ge0\}$. It is nonempty: $k=-|a|$ gives $a+b|a|\ge a+|a|\ge0$.\\
By WOP, $S$ has a least element $r=a-bq$, and $r\ge0$ by definition.\\
Suppose $r\ge b$. Then $r-b=a-b(q+1)\ge0$ lies in $S$ and is smaller than $r$, contradicting minimality. So $r<b$.\\
\emph{Hence $a=bq+r$ with $0\le r<b$.}}
\end{frame}

% ================= 12. Induction & Invariants =================
\begin{frame}{12. Induction and Invariants (Q1--Q2)}
\footnotesize
\qa{1}{Prove $1^3+2^3+\cdots+n^3=\left(\frac{n(n+1)}{2}\right)^2$ for all $n\ge1$.}
{\emph{Base} $n=1$: both sides equal 1.\\
\emph{Step:} assume it holds for $n$. Then $\sum_{i=1}^{n+1}i^3=\frac{n^2(n+1)^2}{4}+(n+1)^3=(n+1)^2\cdot\frac{n^2+4n+4}{4}=\frac{(n+1)^2(n+2)^2}{4}$,\\
which is the formula for $n+1$. \emph{By induction it holds for all $n\ge1$.}}
\qa{2}{Prove $7^n-1$ is divisible by 6 for all $n\ge0$.}
{\emph{Base} $n=0$: $7^0-1=0$, divisible by 6.\\
\emph{Step:} assume $6\mid7^n-1$. Then $7^{n+1}-1=7(7^n-1)+6$. The first term is divisible by 6 by hypothesis and the second is $6$. So $6\mid7^{n+1}-1$.}
\end{frame}

\begin{frame}{12. Induction and Invariants (Q3--Q4)}
\footnotesize
\qa{3}{Prove $2^n>n^2$ for all $n\ge5$.}
{\emph{Base} $n=5$: $32>25$.\\
\emph{Step:} assume $2^n>n^2$ with $n\ge5$. Then $2^{n+1}=2\cdot2^n>2n^2$. It suffices that $2n^2\ge(n+1)^2$, i.e.\ $n^2\ge2n+1$, i.e.\ $(n-1)^2\ge2$, true for $n\ge3$. So $2^{n+1}>(n+1)^2$.}
\qa{4 (invariant)}{A robot starts at $(0,0)$ and moves $(+1,+2)$ or $(+2,+1)$. Can it reach $(1,1)$? $(3,3)$?}
{Each move adds 3 to $x+y$. \emph{Invariant:} after any number of moves, $x+y\equiv0\pmod3$ (true at the start, preserved by each move).\\
$(1,1)$: $x+y=2\not\equiv0$, so it is unreachable.\\
$(3,3)$: $x+y=6\equiv0$, and the moves $(+1,+2)$ then $(+2,+1)$ reach it.}
\end{frame}

\begin{frame}{12. Induction and Invariants (Q5)}
\footnotesize
\qa{5 (invariant)}{The numbers $1,2,\dots,10$ are on a board. Repeatedly erase two numbers $a,b$ and write $|a-b|$. What can you say about the last number?}
{\emph{Invariant:} the parity of the sum of all numbers on the board.\\
Replacing $a,b$ by $|a-b|$ changes the sum by $a+b-|a-b|$, which is $2\min(a,b)$, an even number. So the parity of the sum never changes.\\
The initial sum is $55$, which is odd. The last remaining number is the sum, so it is \textbf{odd}.}
\end{frame}

% ================= 13. Strong Induction =================
\begin{frame}{13. Strong Induction (Q1--Q2)}
\footnotesize
\qa{1}{A shop accepts only 4-rupee and 7-rupee coins. Prove every amount $n\ge18$ can be paid exactly.}
{Claim: $n=4a+7b$ for some $a,b\ge0$.\\
\emph{Base cases:} $18=4+14$, $19=12+7$, $20=20$, $21=21$ (using 4 coins and 7 coins as needed).\\
\emph{Step:} let $n\ge22$ and assume the claim for all $18\le m<n$. Since $n-4\ge18$, $n-4=4a+7b$, so $n=4(a+1)+7b$.\\
Four base cases are needed because the step reaches back by 4.}
\qa{2}{A pile of $n\ge1$ coins is split repeatedly into two nonempty piles until all piles are single coins. Splitting piles of sizes $x$ and $y$ scores $xy$. Prove the total score is $\frac{n(n-1)}{2}$ however the splits are chosen.}
{\emph{Base} $n=1$: no split, score $0=\frac{1\cdot0}{2}$.\\
\emph{Step:} the first split gives piles of $k$ and $n-k$ ($1\le k<n$) and scores $k(n-k)$. By the strong hypothesis the rest score $\frac{k(k-1)}{2}+\frac{(n-k)(n-k-1)}{2}$.\\
Total $=\frac{2k(n-k)+k^2+(n-k)^2-n}{2}=\frac{n^2-n}{2}$.}
\end{frame}

\begin{frame}{13. Strong Induction (Q3--Q4)}
\footnotesize
\qa{3}{With $F_1=F_2=1$ and $F_n=F_{n-1}+F_{n-2}$, prove $F_n<2^n$ for all $n\ge1$.}
{\emph{Base cases:} $F_1=1<2$ and $F_2=1<4$.\\
\emph{Step:} for $n\ge3$ assume $F_k<2^k$ for all $k<n$. Then $F_n=F_{n-1}+F_{n-2}<2^{n-1}+2^{n-2}<2^{n-1}+2^{n-1}=2^n$.\\
Strong induction is needed because $F_n$ depends on two earlier terms.}
\qa{4}{Prove every integer $n\ge2$ is a product of primes.}
{Assume it holds for all $m$ with $2\le m<n$.\\
If $n$ is prime, it is a product of one prime.\\
Otherwise $n=ab$ with $2\le a,b<n$. By the hypothesis both $a$ and $b$ are products of primes, so $n=ab$ is as well.}
\end{frame}

\begin{frame}{13. Strong Induction (Q5)}
\footnotesize
\qa{5}{Let $J_0=0$, $J_1=1$ and $J_n=J_{n-1}+2J_{n-2}$. Prove $J_n=\frac{2^n-(-1)^n}{3}$.}
{\emph{Base cases:} $J_0=\frac{1-1}{3}=0$ and $J_1=\frac{2+1}{3}=1$.\\
\emph{Step:} for $n\ge2$ assume the formula for $n-1$ and $n-2$. Then\\
$3J_n=[2^{n-1}-(-1)^{n-1}]+2[2^{n-2}-(-1)^{n-2}]=2^n+(-1)^n-2(-1)^n=2^n-(-1)^n$,\\
using $-(-1)^{n-1}=(-1)^n$ and $(-1)^{n-2}=(-1)^n$. Dividing by 3 gives the formula.}
\end{frame}

% ================= 14. Structural Induction =================
\begin{frame}{14. Structural Induction (Q1--Q2)}
\footnotesize
\qa{1}{Define string length recursively and prove $|xy|=|x|+|y|$.}
{Define $|\varepsilon|=0$ and $|wa|=|w|+1$ for a symbol $a$. Induct on the structure of $y$.\\
\emph{Base} $y=\varepsilon$: $|x\varepsilon|=|x|=|x|+0=|x|+|\varepsilon|$.\\
\emph{Step} $y=wa$ with $|xw|=|x|+|w|$: $|xwa|=|xw|+1=|x|+|w|+1=|x|+|wa|$.}
\qa{2}{A full binary tree is a single node, or a root with two full subtrees. Prove \#leaves $=$ \#internal nodes $+1$.}
{\emph{Base:} a single node has 1 leaf and 0 internal nodes, and $1=0+1$.\\
\emph{Step:} let $T$ have subtrees $T_1,T_2$ with leaves $\ell_1,\ell_2$ and internal nodes $i_1,i_2$, where $\ell_j=i_j+1$. Then $T$ has $\ell_1+\ell_2=i_1+i_2+2$ leaves and $i_1+i_2+1$ internal nodes (the new root). So leaves $=$ internal $+1$.}
\end{frame}

\begin{frame}{14. Structural Induction (Q3--Q4)}
\footnotesize
\qa{3}{Well-formed formulas are built from variables, $(\neg A)$ and $(A\ast B)$. Prove each wff has equally many ``('' and ``)''.}
{\emph{Base:} a variable has 0 of each.\\
\emph{Step $(\neg A)$:} $A$ has $n$ of each by hypothesis; $(\neg A)$ adds one ``('' and one ``)'', giving $n+1$ of each.\\
\emph{Step $(A\ast B)$:} $A$ has $n$ and $B$ has $m$ of each; adding one pair gives $n+m+1$ of each.}
\qa{4}{Let $S$ be defined by: $3\in S$; if $x,y\in S$ then $x+y\in S$. Prove every element of $S$ is divisible by 3. Which numbers does $S$ contain?}
{\emph{Base:} $3\mid3$. \emph{Step:} if $3\mid x$ and $3\mid y$ then $3\mid x+y$.\\
\emph{Contents:} $S=\{3,6,9,\dots\}$, the positive multiples of 3, since $3k=3+3+\cdots+3$ ($k$ terms) is built by repeated use of the rule.}
\end{frame}

\begin{frame}{14. Structural Induction (Q5)}
\footnotesize
\qa{5}{Define $\mathrm{rev}(\varepsilon)=\varepsilon$ and $\mathrm{rev}(wa)=a\cdot\mathrm{rev}(w)$. Prove $\mathrm{rev}(xy)=\mathrm{rev}(y)\cdot\mathrm{rev}(x)$.}
{Induct on $y$.\\
\emph{Base} $y=\varepsilon$: $\mathrm{rev}(x)=\varepsilon\cdot\mathrm{rev}(x)=\mathrm{rev}(\varepsilon)\cdot\mathrm{rev}(x)$.\\
\emph{Step} $y=wa$, assuming $\mathrm{rev}(xw)=\mathrm{rev}(w)\cdot\mathrm{rev}(x)$:\\
$\mathrm{rev}(xwa)=a\cdot\mathrm{rev}(xw)=a\cdot\mathrm{rev}(w)\cdot\mathrm{rev}(x)=\mathrm{rev}(wa)\cdot\mathrm{rev}(x)$.}
\end{frame}

% ================= Closing =================
\begin{frame}
  \centering
  \vspace{1cm}
  {\Huge Thank You}\\[1cm]
\end{frame}

\end{document}